\subsection{A precise obstruction to deleting native stack traversal}
\label{sec:tail-deletion}

The truth-axiom elimination does not make the remaining presentation an
arbitrary deletion target. In particular, remove the native tail equation
\[
 V(P(x,s),A(j,d))=V(s,j)
 \tag{Tail}
\]
from $\mathcal E_{688}$, and call the resulting six-equation theory
$\mathcal E^-$. There is an explicit separating $\mathcal E^-$ model in
ZFC without any consistency assumption. Thus this second deletion would
destroy the intended consistency calibration, assuming ZFC is consistent.

\paragraph{The carrier and Boolean operations.}
Use the carrier $\mathbb N$. Code a finite set by its binary bit mask $m$;
the even number $2m$ represents its complement in $\mathbb N$, while
$2m+1$ represents the finite set itself. This is the finite/cofinite
Boolean algebra. Its top is the native value $d=0$, its bottom is $1$,
complement toggles the last bit, and $N$ is NAND. Sh1 therefore holds and
$N(d,d)=1\ne0=d$.

Let $\pi(x,y)=\frac{(x+y)(x+y+1)}2+y$ and define
\[
 P(x,s)=3\pi(x,s)+1,\qquad A(i,j)=3\pi(i,j)+2,
 \qquad Q(a)=3a+3.
\]
These three constructor ranges partition the positive naturals and are
disjoint; each constructor code exceeds its arguments. Interpret
$V(s,i)$ as the first decoded component when $s$ is a $P$ code and as
$0$ otherwise, \emph{independently of $i$}. The native head equation holds.

Define $\tau(a)\in\{0,1\}$ by recursion on the natural class code:
\[
 \tau(0)=1,\quad
 \tau(P(a,b))=N(\tau(a),\tau(b)),\quad
 \tau(A(i,j))=1,\quad
 \tau(Q(a))=N(\tau(a),\tau(a)).
\]
The recursion is well founded because its children are smaller.
Set $p_\varepsilon(s,a)=\tau(a)$, so incidence is independent of the stream.
Set $p_E(x,y)=0$ when $x<y$ and $p_E(x,y)=1$ otherwise; the Boolean value
$0$ is truth in this coding. Finally let
\[
 W(s,a)=0,\qquad\mathsf{Target}(b)=0,\qquad
 \mathsf{Next}(b,z)=z+2,\qquad\mathsf{Mem}(b,z)=z+1.
\]

\paragraph{Verification for all arguments.}
The NAND-class equation holds by the definition of $\tau$. For the atom
equation, $V(s,i)=V(s,j)$, so the right side is the false value $1$;
the left side is $\tau(A(i,j))=1$. In the complemented-binder equation,
the two selected and arbitrary evaluations are both $\tau(a)$, hence its
right side is $N(\tau(a),\tau(a))=\tau(Q(a))$, exactly its left side.

For a fixed descriptor $b$, the relation value
\[
 R_b(x,y)=p_\varepsilon(P(y,P(x,b)),V(b,d))
          =\tau(V(b,d))
\]
is constant in both $x$ and $y$. If that constant is truth, the bad branch
of the closure equation is true at every $z$: its three conjuncts are
the true matrix value, $z+1<z+2$, and the negation of $z+1<z$.
If the constant is false, the target $h=0$ has no members, so both closure
implications and its unrestricted support clause are true. In either
case the retained factored closure equation holds. This argument covers
every descriptor and every natural assignment, including noncanonical
selectors and streams.

Nevertheless $(\mathrm{Tail})$ fails at $x=1$, $s=0$, $j=0$: its left
side is $1$, while its right side is $0$. All six retained equations hold
and the original ground endpoints are separated. The model is explicitly
defined using arithmetic, finite/cofinite Boolean operations, and a
well-founded recursion, so ZFC proves its existence directly. Consequently
ZFC proves $\mathcal E^-\not\vdash N(d,d)=d$. If this were also provably
equivalent in ZFC to $\operatorname{Con}(\mathrm{ZFC})$, ZFC would prove
its own consistency.

The accompanying checker evaluates the six literal AST equations on
$6{,}208$ assignments with all variables in $\{0,\ldots,7\}$ and verifies
the displayed tail failure and endpoint separation. Those finite cases
check the implementation against the stated construction; the preceding
all-natural argument proves the model theorem. This control identifies
the logical role of actual coordinate traversal and explains why a
seemingly small deletion cannot be priced as another valid improvement.
